% ============================================================
% CHAPTER 5: SRE METHODOLOGY & REQUIREMENTS
% MCDA 5585 / 5586 Major Project Report
% Bhavik Kantilal Bhagat | A00494758 | Saint Mary's University
% ============================================================

\chapter{Methodologies \& Requirements}
\label{chap:methodologies}

This chapter documents the engineering methodology that underpinned my
internship. It begins with the Site Reliability Engineering principles
I adopted from the team's knowledge-transfer sessions and explains how
I operationalised those principles for the Citco Intelligent
Automation platform. It then characterises the observability gap that
existed at the start of my engagement, identifies the stakeholders who
shaped the requirements, describes how I elicited requirements, and
closes with the complete functional and non-functional requirements that
governed the build.

% ────────────────────────────────────────────────────────────
\section{SRE Principles Applied}
\label{sec:sre_principles}
% ────────────────────────────────────────────────────────────

Site Reliability Engineering is an operational discipline that applies
software engineering thinking to the problem of keeping production
systems reliable at scale. Its conceptual vocabulary --- golden signals,
Service Level Indicators (SLIs), Service Level Objectives (SLOs), and
error budgets --- provided the theoretical backbone for every design
decision I made during this internship.

\subsection{The Four Golden Signals}

The four golden signals were introduced to the team by mentor Kristia
Marie Labos (Kri) during a dedicated knowledge-transfer session on
\textbf{May~25, 2026}. They define the minimum instrumentation required
to characterise the health of any service, regardless of the underlying
technology.

\begin{enumerate}
    \item \textbf{Latency} --- \emph{How long does a request take?} For
          the IA platform, latency manifests as Lambda function duration
          (measured at the P50, P95, and P99 percentiles), and as ECS
          service response time for API-facing containers. Sustained
          increases in P99 latency are an early warning of resource
          saturation or dependency degradation before errors appear.

    \item \textbf{Traffic} --- \emph{How much demand is on the system?}
          Traffic is measured as Lambda invocation count per unit time,
          SQS message ingestion rate for queue-backed automations, and
          MSK throughput (bytes in/out per second) for the Meridian
          Kafka pipeline. Traffic anomalies --- sudden spikes or
          unexpected drops --- often precede error and saturation
          signals.

    \item \textbf{Errors} --- \emph{What is the rate of failed
          requests?} The IA platform tracks errors as Lambda error and
          throttle counts, ECS task failure counts, and
          \texttt{NumberOfMessagesFailed} on SQS queues. An error rate
          that crosses a defined threshold indicates a code or
          infrastructure defect requiring immediate attention.

    \item \textbf{Saturation} --- \emph{How full is the service?}
          Saturation is the signal most relevant to the Apr~17,~2026
          Meridian incident. It is measured as Lambda concurrent
          execution count relative to the account concurrency limit,
          ECS task CPU and memory utilisation, MSK consumer lag (the
          number of messages awaiting processing on a Kafka topic), and
          SQS queue depth with oldest-message age. A service approaching
          saturation degrades before it fails --- making this the most
          actionable of the four signals.
\end{enumerate}

\subsection{Service Level Indicators and Objectives}

A \textbf{Service Level Indicator (SLI)} is a carefully defined
quantitative measure of some aspect of the level of service being
delivered. An SLI is not a raw metric --- it is the specific metric
chosen to represent one of the four golden signals for a particular
service. A \textbf{Service Level Objective (SLO)} is a target value
or range of values for an SLI, measured over a rolling time window.
Together, SLIs and SLOs transform vague aspirations such as
``the system should be reliable'' into falsifiable engineering
commitments.

For the IA platform, I adopted the following SLOs as the design
targets for the observability build:

\begin{itemize}
    \item \textbf{Lambda error rate SLO}: Fewer than 1\% of Lambda
          invocations must result in an error over any 24-hour rolling
          window. This translates directly to the CloudWatch alarm
          threshold applied to the \texttt{Errors} metric with
          \texttt{ErrorRate = Errors / Invocations}.

    \item \textbf{MSK consumer lag SLO}: The Kafka consumer lag for
          any Meridian consumer group must remain below
          \textbf{10,000~messages} at any point in time. This threshold
          was derived from the Apr~17 post-mortem: the system reached
          9~million messages in backlog before the business team
          escalated, and a threshold of 10,000 would have triggered an
          alarm within minutes of the first deviation.

    \item \textbf{ECS CPU utilisation SLO}: No ECS service should
          sustain CPU utilisation above \textbf{80\%} over a 5-minute
          evaluation period. Breaching this threshold activates the
          auto-scaling policy and triggers an SNS notification to the
          SRE team.

    \item \textbf{SQS queue depth SLO}: No production SQS queue should
          accumulate more than a configurable maximum of unprocessed
          messages before an alarm fires. The threshold is set per-queue
          based on the expected processing throughput for that
          automation.
\end{itemize}

\subsection{Dynatrace Parity Analysis}
\label{sec:dynatrace_parity}

A key framing question for my internship --- introduced in
Chapter~\ref{chap:executive_summary} and developed here in methodological
detail --- is how closely an AWS-native observability stack can replicate
the capabilities of Dynatrace, and at what incremental cost. The PoC
investigation evaluated 25 Dynatrace features against five AWS tool
combinations to quantify achievable parity.

Table~\ref{tab:dynatrace_parity} summarises the parity analysis for the
ten most operationally relevant Dynatrace capabilities to the IA platform.

\begin{table}[ht]
\centering
\caption{Dynatrace Capability Parity --- AWS-Native Stack vs Dynatrace (selected features)}
\label{tab:dynatrace_parity}
\begin{tabular}{p{4.5cm}p{1.4cm}p{1.4cm}p{1.4cm}p{1.8cm}p{1.4cm}}
\toprule
\textbf{Dynatrace Feature} & \textbf{CW} & \textbf{+Grafana} & \textbf{+X-Ray} & \textbf{+App Signals} & \textbf{+Tempo} \\
\midrule
Alerting                  & \checkmark & \checkmark & --- & --- & --- \\
Infrastructure Metrics    & \checkmark & \checkmark & --- & --- & --- \\
Log Analytics             & \checkmark & \checkmark & --- & --- & --- \\
Container Monitoring      & \checkmark & \checkmark & --- & \checkmark & --- \\
Error Rate Tracking       & \checkmark & \checkmark & \checkmark & \checkmark & --- \\
Latency Percentiles       & --- & \checkmark & \checkmark & \checkmark & \checkmark \\
Distributed Trace Waterfall & --- & --- & \checkmark & --- & \checkmark \\
Service Map / Topology    & --- & --- & \checkmark & \checkmark & --- \\
SLO Management            & --- & --- & --- & \checkmark & --- \\
Dependency Mapping        & --- & --- & \checkmark & \checkmark & --- \\
\midrule
\textbf{Parity (\% of 25 features)} & 30\% & 33\% & 60\% & 80\% & 83\% \\
\textbf{Monthly Cost (est.)} & \$49 & \$68 & \$77 & \$232 & \$269 \\
\bottomrule
\end{tabular}
\end{table}

\noindent
Three capabilities remain \textbf{permanently outside the AWS-native reach}:
Davis~AI causal analysis (proprietary ML, no AWS equivalent); automatic
baselining (Dynatrace learns per-service baselines over weeks; AWS alarms
require manual thresholds); and session replay (full DOM capture, not
available in any OSS tool). These account for approximately 9--10\% of
the functional gap that no combination of AWS tools can close.

\paragraph{Adopted stack.}
The PoC I delivered in this internship corresponds to the
\textbf{CloudWatch + Grafana + X-Ray} column: 60\% parity at approximately
\$77/month. Application Signals (the path to 80\% parity) requires deploying
an ADOT sidecar on each ECS service and is scoped as post-internship work.
The gap-to-tooling mapping is used throughout Chapter~\ref{chap:implementation}
and Chapter~\ref{chap:achievements} to contextualise what the PoC does and does
not cover.

\subsection{Error Budgets}

An \textbf{error budget} is the complement of the SLO: it is the
maximum permissible unreliability for a service over a given window.
If the SLO is that Lambda error rate must be below 1\% over 24~hours,
then the error budget is the 1\% of invocations that are permitted to
fail before the SLO is breached. When the error budget is consumed ---
either by a single outage or by sustained degradation --- the SRE
engineering principle dictates that all non-reliability work (feature
development, infrastructure experiments) should be paused in favour of
reliability remediation.

In the context of my IA internship, I applied error budgets
\emph{conceptually} rather than through automated budget-burn tracking.
The Grafana dashboard exposes the current SLI values for each monitored
service, and the CloudWatch alarms fire when the system approaches the
SLO threshold --- effectively providing a manual signal that the error
budget is being drawn down. Full automated budget-burn rate alerting
(which would require custom CloudWatch Metric Math expressions computing
burn rate against a rolling budget) was identified as a post-internship
enhancement.

% ────────────────────────────────────────────────────────────
\section{Observability Gap Analysis}
\label{sec:observability_gap}
% ────────────────────────────────────────────────────────────

A prerequisite for building any monitoring platform is understanding the
current state of observability. Before I wrote a single line of code
for IISS-487, I performed a structured gap analysis to characterise
what the IA platform could and could not see about itself.

\subsection{Pre-Internship Monitoring Posture}

At the commencement of my internship in Mar~2026, the IA-SRE team's
monitoring posture exhibited the following deficiencies:

\begin{itemize}
    \item \textbf{No CloudWatch alarms on any Meridian service.} The
          Meridian Kafka-backed event pipeline --- the most
          business-critical component of the IA platform --- had zero
          proactive alerting. No alarms existed on MSK consumer lag,
          ECS task failure counts, Lambda error rates, or SQS queue
          depth for any Meridian-owned resource.

    \item \textbf{No centralized dashboard.} Monitoring was entirely
          fragmented. Operators who wished to check the health of a
          service were required to navigate to that service's individual
          CloudWatch console page. There was no single pane of glass
          from which to assess the overall health of the platform.

    \item \textbf{Incomplete Container Insights coverage.} Of the six
          Meridian ECS clusters, only one
          (\texttt{meridian-compute-cluster}, enabled since Jul~2024)
          had Amazon CloudWatch Container Insights activated. The
          remaining five clusters were invisible to task-level CPU and
          memory telemetry.

    \item \textbf{No MSK consumer lag monitoring.} Despite Kafka
          consumer lag being the primary leading indicator of Meridian
          processing health --- as demonstrated conclusively by the
          Apr~17 incident --- no metric or alarm existed for this
          signal. The \texttt{EstimatedMaxTimeLag} and
          \texttt{SumOffsetLag} metrics available natively from Amazon
          MSK were not surfaced anywhere.

    \item \textbf{No SLO tracking or error budget calculation.}
          The team had no defined SLIs, no published SLO targets, and
          therefore no basis for calculating error budgets or
          determining when the platform was in an unacceptable
          reliability state.

    \item \textbf{No automated resource inventory.} The inventory of
          Lambda functions, ECS clusters, and EC2 instances was
          maintained manually via Confluence pages and spreadsheets.
          This made it impossible to know, at any given moment, exactly
          which resources existed across all six platforms.

    \item \textbf{RPA bot monitoring limited to orchestrator UIs.}
          Blue Prism and UiPath bots were monitored exclusively through
          their respective orchestrator platforms (Blue Prism Orchestrator and
          UiPath Cloud). No CloudWatch integration existed, and there
          was no mechanism for the SRE team to receive proactive
          notification of bot failures or queue backlogs outside the
          orchestrator UI.

    \item \textbf{Reactive operational posture.} The cumulative effect
          of the above gaps was that the team's primary incident
          detection mechanism was business-team escalation. When a fund
          processing automation failed silently, the first signal
          received by the SRE team was a ticket or message from
          operations staff in Dublin, London, or Halifax who had noticed
          anomalous fund processing outputs. This reactive posture
          introduced structural latency between failure onset and
          engineering response.
\end{itemize}

\subsection{Gap Analysis Summary}

Table~\ref{tab:gap_analysis} summarises the pre-internship gaps and the
post-internship resolutions for each monitoring domain.

\begin{table}[ht]
\centering
\caption{Observability Gap Analysis --- Pre- vs.\ Post-Internship}
\label{tab:gap_analysis}
\begin{tabular}{p{3.5cm}p{4.5cm}p{4.5cm}}
\toprule
\textbf{Domain} & \textbf{Pre-Internship State} & \textbf{Post-Internship State} \\
\midrule
CloudWatch Alarms & 0 alarms on any IA resource & 38 alarms across 7 AWS service categories \\
\midrule
Centralized Dashboard & None & 14~rows, 110+~panels (Grafana AMG) \\
\midrule
Container Insights & 1 of 6 Meridian clusters & 5 of 6 Meridian clusters \\
\midrule
MSK Consumer Lag & No monitoring & CloudWatch alarm on \texttt{EstimatedMaxTimeLag} ($<$30~sec detection) \\
\midrule
SLO Tracking & None & SLI/SLO targets defined; dashboard exposes current values \\
\midrule
Resource Inventory & Manual (Confluence) & Automated REST API (Lambda, ECS, EC2 across all platforms) \\
\midrule
Incident Detection & Business team escalation & Proactive alarm with SNS routing \\
\bottomrule
\end{tabular}
\end{table}

% ────────────────────────────────────────────────────────────
\section{Stakeholder Identification}
\label{sec:stakeholders}
% ────────────────────────────────────────────────────────────

Requirements for an internal SRE tooling platform are shaped by a
relatively small but technically diverse set of stakeholders.
Unlike external-facing software products, the IA observability
platform has no end-user customer outside Citco. All stakeholders
are internal --- either direct members of the IA-SRE team or
operational teams who depend on the automations that the platform
monitors. Table~\ref{tab:stakeholders} enumerates the key
stakeholders, their roles, and their specific interests in the
project.

\begin{table}[ht]
\centering
\caption{Stakeholder Identification}
\label{tab:stakeholders}
\begin{tabular}{p{3.8cm}p{3.0cm}p{5.5cm}}
\toprule
\textbf{Stakeholder} & \textbf{Role} & \textbf{Interest / Influence} \\
\midrule
Kishor Deotale & Manager, CTM; Industry Supervisor & Primary sponsor; sets engineering priorities; escalated IISS-487 to Critical on May~6, 2026; final acceptance authority \\
\midrule
Ho Hoi Leung (Horace) & RPA Platform Lead & RPA platform subject matter expert; provided KT on Blue Prism and UiPath architectures; defined RPA monitoring scope \\
\midrule
Kristia Marie Labos (Kri) & SRE Methodology Lead & Provided golden signals and SLI/SLO training (May~25 session); reviewed dashboard design; validated alarm threshold selection \\
\midrule
IA Operations Team & Internal SRE consumers & Primary users of the Grafana dashboard and alarm notifications; operational first responders to CloudWatch alarms \\
\midrule
Business Operations Teams (Dublin, London, Halifax) & Fund admin automation consumers & Downstream consumers of the automations; surface incidents under the reactive model; motivated the shift to proactive monitoring \\
\midrule
Meridian Platform Team & Kafka/ECS architecture owners & Provided architecture context for MSK consumer group structure, ECS cluster topology, and OpenSearch integration \\
\midrule
Co-Interns (Sridhar Vadla, Soundarya Venkataraman) & Collaborative development partners & Parallel development of related monitoring features; design review and test validation \\
\bottomrule
\end{tabular}
\end{table}

% ────────────────────────────────────────────────────────────
\section{Requirements Elicitation}
\label{sec:requirements_elicitation}
% ────────────────────────────────────────────────────────────

All stakeholders for this project were internal to Citco. This
constrained the elicitation methods to those appropriate for an
in-house engineering context: structured knowledge-transfer sessions,
incident post-mortems, and JIRA story refinement, rather than formal
user interviews or focus groups.

\subsection{Elicitation Methods}

\paragraph{Knowledge-Transfer Sessions (Training Week, Apr 13--18, 2026).}
The first week of full-time engagement was structured as an intensive
knowledge-transfer programme covering the full breadth of the IA
platform. Sessions with Horace on the Blue Prism and UiPath
architectures, with Bhanu Teja Murari on CitcoWorks and the Event Store
platform, and with the Meridian team on the Kafka event pipeline
(comprising the Event Routing Service, Action Processor, and OpenSearch
consumer) provided the architectural grounding needed to define
meaningful SLIs. Key requirements insight from training week: the
\AE xeo Treasury platform (cloud-based wire approval and fund movement)
and the MESO platform (high-memory autoscale issue, IISS-538) each
had distinct resource profiles requiring separate monitoring strategies.

\paragraph{Incident Post-Mortem (Apr 17, 2026).}
The Meridian production incident yielded the single most important
requirements insight of the entire engagement: \textbf{MSK consumer
lag is the primary leading indicator of Meridian health}. The
post-mortem made concrete the latency between signal onset and
detection under the reactive model, and directly specified the target
for the lag alarm evaluation period (NFR2: $\leq$1~minute).

\paragraph{Golden Signals Session (May 25, 2026).}
Kri's dedicated SRE methodology session translated the theoretical
four golden signals into concrete metric choices for the IA platform.
This session produced the initial SLI/SLO target table and established
the design constraint that all metrics must be queryable through
Amazon CloudWatch --- ruling out external time-series databases such as
Prometheus and anchoring the architecture to CloudWatch Metric Insights
as the sole data source.

\paragraph{JIRA Story Refinement.}
IISS-487 went through multiple refinement cycles between May~6 and
May~11 (infrastructure deployment date). These sessions with Kishor
and Kri surfaced the infrastructure-as-code requirement (FR7), the
multi-environment reuse requirement (NFR5), and the credential
management constraint (NFR7). The elevation of IISS-487 to Critical
priority on May~6 also functionally set NFR1 (dashboard refresh
latency $\leq$30~seconds) as a hard constraint rather than an aspirational
target.

\subsection{Key Insights}

Three insights from the elicitation process materially shaped the
resulting design:

\begin{enumerate}
    \item \textbf{Tag-based resource discovery is non-trivial at Citco
          scale.} The training week revealed that the \texttt{CITCOWORKS}
          and \texttt{MERIDIAN} environments share the
          \texttt{``Shared Cloud Artifacts''} BudgetCode tag, meaning
          a simple \texttt{BudgetCode} lookup conflates resources from
          two operationally distinct environments. This directly
          motivated FR9 (tag-based discovery with \texttt{AppName}
          fallback).

    \item \textbf{The primary audience for the dashboard is the SRE
          team, not business operators.} Business operations teams would
          continue to use their existing domain-specific tools; the
          dashboard needed to be useful for the SRE first-responder,
          not designed as a business reporting tool. This shaped the
          panel selection (technical metrics over business KPIs) and
          the deep-link design (panels link directly to the relevant AWS
          Console page).

    \item \textbf{Infrastructure-as-Code is a first-class requirement,
          not an after-thought.} The team had experienced the cost of
          manually provisioned monitoring infrastructure in previous
          projects (configuration drift, undocumented alarms). The JIRA
          refinement sessions made FR7 an explicit acceptance criterion
          rather than a ``nice to have.''
\end{enumerate}

% ────────────────────────────────────────────────────────────
\section{Functional and Non-Functional Requirements}
\label{sec:requirements}
% ────────────────────────────────────────────────────────────

The following requirements were derived from the elicitation activities
described above. They governed the design and acceptance criteria for
all three work streams of my internship.

\subsection{Functional Requirements}

Table~\ref{tab:functional_requirements} enumerates the nine functional
requirements, each with the primary deliverable that satisfies it.

\begin{table}[ht]
\centering
\caption{Functional Requirements}
\label{tab:functional_requirements}
\begin{tabular}{p{0.8cm}p{6.8cm}p{4.5cm}}
\toprule
\textbf{ID} & \textbf{Requirement} & \textbf{Satisfied By} \\
\midrule
FR1 & Dashboard must display real-time metrics for Lambda functions (invocations, errors, duration, throttles, concurrent executions) per budget code & Grafana AMG dashboard, Lambda rows (IISS-487) \\
\midrule
FR2 & Dashboard must display ECS service CPU and memory utilisation per cluster and service & Container Insights + ECS rows in Grafana (IISS-487, IISS-483) \\
\midrule
FR3 & Dashboard must display MSK consumer lag per consumer group and topic partition & MSK rows in Grafana; CloudWatch \texttt{EstimatedMaxTimeLag} alarm (IISS-487) \\
\midrule
FR4 & Dashboard must display SQS queue depth and oldest message age per queue & SQS rows in Grafana; SQS monitoring API (IISS-482) \\
\midrule
FR5 & CloudWatch alarms must be provisioned for all monitored resource types with SNS notification routing to SRE team & 38 CloudWatch alarms deployed via CloudFormation (IISS-487) \\
\midrule
FR6 & Dashboard must support variable-based filtering by budget code, environment, and resource name & Grafana template variables (IISS-487) \\
\midrule
FR7 & All monitoring infrastructure must be deployed as Infrastructure-as-Code via CloudFormation & 4 CloudFormation stacks; Lambda-backed IaC provisioner; S3-staged templates (IISS-487) \\
\midrule
FR8 & Inventory API must enumerate all Lambda, ECS, EC2, and S3 resources per budget code via REST endpoints & Flask REST API, discoverer registry (IISS-507) \\
\midrule
FR9 & Inventory API must support tag-based resource discovery with \texttt{AppName} fallback for shared BudgetCode environments & Two-stage filter logic in \texttt{CITCOWORKS}/\texttt{MERIDIAN} discoverers (IISS-507) \\
\bottomrule
\end{tabular}
\end{table}

\subsection{Non-Functional Requirements}

Table~\ref{tab:nonfunctional_requirements} enumerates the seven
non-functional requirements, each with the measurable target and the
achieved value where applicable.

\begin{table}[ht]
\centering
\caption{Non-Functional Requirements}
\label{tab:nonfunctional_requirements}
\begin{tabular}{p{0.8cm}p{4.5cm}p{2.8cm}p{3.5cm}}
\toprule
\textbf{ID} & \textbf{Requirement} & \textbf{Target Metric} & \textbf{Achieved} \\
\midrule
NFR1 & Dashboard refresh latency must enable sub-minute incident detection & $\leq$30~seconds & \textcolor{smumaroon}{$\leq$30~sec} (CloudWatch native refresh) \\
\midrule
NFR2 & CloudWatch alarm evaluation period for critical metrics (MSK consumer lag, Lambda error rate) & $\leq$1~minute & \textcolor{smumaroon}{1-minute period} (deployed) \\
\midrule
NFR3 & Inventory API response time for full resource enumeration & $\leq$10~seconds & \textcolor{smumaroon}{$<$10~sec} (async Flask, async AWS SDK) \\
\midrule
NFR4 & Inventory API automated test coverage & $\geq$95\% & \textcolor{smumaroon}{99\%} (562 tests) \\
\midrule
NFR5 & All CloudFormation templates must be parameterized for multi-environment reuse without code changes & Parameterized YAML & \textcolor{smumaroon}{Achieved} (env, budget code, thresholds as parameters) \\
\midrule
NFR6 & Dashboard must remain functional after Grafana version upgrades without manual panel rework & Standard plugin queries only & \textcolor{smumaroon}{Achieved} (survived 10.4$\to$12.4 upgrade with query migration) \\
\midrule
NFR7 & Zero hardcoded credentials --- all secrets via AWS Secrets Manager & No plaintext secrets in code or templates & \textcolor{smumaroon}{Achieved} (all credentials via Secrets Manager; validated in IISS-706 hotfix) \\
\bottomrule
\end{tabular}
\end{table}

\subsection{Requirements Traceability}

All nine functional requirements and all seven non-functional
requirements were satisfied by my delivered work streams. The
test suite for IISS-507 provides quantitative evidence for NFR3
and NFR4. NFR6 was validated in practice when the Grafana workspace
was upgraded from version 10.4 to version 12.4 on May~13, 2026 ---
a major version migration that required updating CloudWatch Metric
Insights query syntax across all panels but did not require any
architectural change to the underlying monitoring design.

The most technically demanding requirement was FR9, which exposed a
fundamental ambiguity in the AWS tagging strategy for the IA platform.
My discovery that \texttt{CITCOWORKS} and \texttt{MERIDIAN} share the
same \texttt{BudgetCode} tag value (\texttt{``Shared Cloud Artifacts''})
was not documented anywhere before I built the inventory API. The
solution --- designating \texttt{AppName} as a mandatory secondary
discriminator for these two environments --- became the standard pattern
for all resource discovery going forward and was documented in the
inventory API's \texttt{README} as a platform-wide architectural note.

NFR7 received an unexpected real-world validation during my
IISS-706 hotfix on Jun~8, 2026. The CAIS Deadline incident revealed
that the ECS container's credential loading pattern (reading from
Secrets Manager once at module import time via \texttt{constants.py})
was fragile under secret rotation. The hotfix I delivered introduced fail-fast
credential validation --- a design improvement that reinforced NFR7
by ensuring that absent or expired credentials from Secrets Manager
produce an immediate, visible failure rather than a silent downstream
authentication error.

\subsection{Out-of-Scope Requirements}

Two requirements categories were explicitly deferred as post-internship
work:

\begin{itemize}
    \item \textbf{Automated error budget burn-rate alerting.} Computing
          burn rate requires Metric Math expressions that divide the
          rolling error count by the rolling SLO budget. While the SLO
          targets and SLI metrics are now in place, configuring
          burn-rate alarms requires additional CloudFormation work that
          I descoped from IISS-487.

    \item \textbf{RPA bot CloudWatch integration.} Bridging the Blue
          Prism and UiPath orchestrators to CloudWatch (to surface
          bot queue depth and failure counts as native AWS metrics)
          requires a polling Lambda function or a webhook integration
          with the orchestrator APIs. I identified this during the
          RPA support KT sessions as a high-value enhancement but did
          not include it in the current scope.
\end{itemize}

% ────────────────────────────────────────────────────────────
\section{Limitations and Mitigations}
\label{sec:limitations}
% ────────────────────────────────────────────────────────────

Several CloudWatch and AMG constraints shaped the final architecture and
required deliberate design decisions to work around.

\textbf{CloudWatch \texttt{SEARCH()} visibility gap.} The \texttt{SEARCH()}
function only returns metrics for AWS services with recent activity. Lambda
functions not invoked within the lookback window disappear from the
dashboard entirely --- the same root cause that made the CAIS pipeline
invisible in IISS-430. \textit{Mitigation:} I replaced all
\texttt{SEARCH()} expressions with CloudWatch Metric Insights SQL
(\texttt{SELECT ... GROUP BY FunctionName}), which enumerates all
deployed AWS services in the namespace regardless of recent activity.

\textbf{Metric Insights SQL cannot compute percentiles.} The
\texttt{GROUP BY} mode of Metric Insights only supports \texttt{AVG},
\texttt{SUM}, \texttt{MIN}, and \texttt{MAX}. P50/P99 latency statistics
are unavailable in this mode. \textit{Mitigation:} Latency percentile
panels use CloudWatch standard metrics mode with a
\texttt{FunctionName:~["*"]} wildcard dimension, which restores true
percentiles but reintroduces the visibility gap for idle AWS services ---
an accepted tradeoff documented in the dashboard design.

\textbf{CloudWatch subscription filters blocked by IAM boundary.}
A proposed central log-aggregation design would require 106 metric
subscription filters across the six platforms. Analysis confirmed this
was technically feasible but blocked by the
\texttt{logs:PutSubscriptionFilter} action not being granted within the
team's scoped permission set. \textit{Mitigation:} I adopted CloudWatch
Logs Insights widgets embedded in Grafana panels for log-based error
surfacing, avoiding subscription filters entirely.

\textbf{Secrets Manager 64\,KB payload limit.} The Lambda-backed
CloudFormation provisioner stores alarm templates as a Secrets Manager
payload. As the monitoring scope expanded to cover all AWS service
categories, the generated YAML exceeded the hard 64\,KB limit.
\textit{Mitigation:} I migrated all templates to an S3 bucket with
object versioning; the provisioner fetches templates via
\texttt{s3:GetObject} at deploy time.

\textbf{Shared \texttt{BudgetCode} tag across platforms.} The
\texttt{meridian} and \texttt{citcoworks} platforms share the internal
tag value \texttt{"Shared Cloud Artifacts"}, making a single-stage tag
lookup ambiguous when discovering AWS services in those platforms.
\textit{Mitigation:} I designated \texttt{AppName} as the mandatory
secondary discriminator for all AWS service discovery in those platforms,
documented as a platform-wide governance constraint in the inventory API
\texttt{README}.
